\documentclass{rptuma}
% Shared preamble for the three structure proposals. Built against ../rptuma.cls by build.sh,
% so each proposal renders exactly as a thesis written in the department template would.
\usepackage{xcolor}
\usepackage{booktabs}
\usepackage{tabularx}
\usepackage{enumitem}

% The class loads its logo as img/RPTUlogo; the image lives one directory up.
\graphicspath{{../}}

\definecolor{materialgray}{gray}{0.38}
\newcommand{\material}[1]{\par{\raggedright\small\sffamily\color{materialgray}\textbf{Material:} #1\par}}
\newcommand{\repo}[1]{\path{#1}}
\newcommand{\us}{\,\textmu s}

\setlist{nosep,leftmargin=*}
% Repository paths and identifiers are long unbreakable words; loose spacing beats margin overruns.
\AtBeginDocument{\sloppy}
\setcounter{tocdepth}{2}
% KOMA appends a period to section numbers once appendices exist; the department sample has none.
\KOMAoptions{numbers=noendperiod}
\addto\captionsenglish{\renewcommand{\bibname}{Sources Used for This Proposal}}

\newdate{thesisDate}{07}{10}{2026}
\author{Author's Name}
\department{University of Kaiserslautern-Landau\\Department of Computer Science\\Distributed Systems Lab}
\supervisortitle{Examiner:}
\supervisor{Prof. Dr.-Ing. Jens B. Schmitt}
\secondsupervisor{M. Sc. Supervisor Name}

\newcommand{\frontmatternote}{%
\section*{Front matter (unchanged from the department template)}
\begin{itemize}
  \item \textbf{Title page} and \textbf{Eigenst\"andigkeitserkl\"arung}, as generated by \texttt{rptuma.cls}.
        The declaration requires disclosing generative-AI use beyond grammar and style.
  \item \textbf{Abstract} (English): the problem (deception at the web edge has an unknown cost), the
        approach (one mechanism on four edges, levelled to identical behaviour), the method (paired
        benchmark on three measurement planes) and the headline numbers.
  \item \textbf{Zusammenfassung}: the same abstract in German.
  \item \textbf{Contents}, \textbf{List of Figures}, \textbf{List of Tables}.
  \item \textbf{Glossary} (optional, the class's \texttt{glossar} environment): WADM, honeytoken, edge,
        tier, VU, iteration, SLO, cgroup, LuaJIT, PUC Lua, proxy-wasm.
\end{itemize}
In the department template's own example, websites are cited through an archived copy
(\texttt{web.archive.org}); do the same for any web source in the list at the end.}


\subject{Master Thesis -- Final Structure}
\title{Final Template: Design, Implementation, Evaluation with Research Questions}

\begin{document}

\maketitle
\pagenumbering{roman}
\setcounter{page}{1}

\chapter*{About This Proposal}

\textbf{Organising idea.} Template~A's chapters unchanged: the thesis tells the story of a
system, from problem through design and implementation on four edges to how well it works. From
Template~B it takes the research questions as the backbone of the evaluation. They are stated in
the introduction, every metric in the methodology is tied to one, the Results chapter has one
section per question ending with an explicit answer, and the conclusion answers each in turn.
Modelled on general thesis guidance \cite{intille_thesis,unibz_guidelines}, systems measurement
papers such as MeshInsight \cite{zhu2022dissecting}, and the RQ convention of empirical studies
\cite{wohlin2012experimentation}.

\textbf{What changed from A.} Section~\ref{sec:AB.rq} states the questions in full;
Section~\ref{sec:AB.metrics} maps metrics to them; Chapter~\ref{chap:AB.results} is regrouped into
RQ1--RQ4; Discussion opens with a synthesis across the questions; Conclusion answers each one.

\begin{center}
\small
\begin{tabularx}{\textwidth}{@{}rXr@{}}
\toprule
Ch. & Title & approx.\ pages \\
\midrule
1 & Introduction & 5 \\
2 & Background & 10 \\
3 & Related Work & 6 \\
4 & Design of WADM & 12 \\
5 & Implementation on Four Edges & 14 \\
6 & Evaluation Methodology & 10 \\
7 & Results (RQ1--RQ4) & 20 \\
8 & Discussion & 6 \\
9 & Conclusion and Future Work & 3 \\
A--D & Appendices & as needed \\
\bottomrule
\end{tabularx}
\end{center}
{\footnotesize Page counts are a rough budget, not a department rule. The unibz guidelines
advise keeping a chapter under about 20 pages and splitting it above 30 \cite{unibz_guidelines}.}

\textbf{How to read the rest.} The table of contents is the proposed structure. Under each
heading is a short note on what to write there, and a \emph{Material} line pointing to the
place in the repository that already holds it (\repo{results/} and \repo{plots/} are short for
\repo{benchmarks/results/} and \repo{benchmarks/results/plots/}).

\frontmatternote

\tableofcontents
\cleardoublepage
\pagenumbering{arabic}

%======================================================================
\chapter{Introduction}
Why deception at the web edge is worth studying, and what this thesis contributes.

\section{Motivation}
Attackers inspect a web application's HTML, cookies and hidden paths before attacking it.
Honeytokens planted there raise high-signal alerts, because no legitimate user touches them. An
edge proxy can plant them without changing the application.

\section{Problem Statement}
Nobody has measured what this costs. Which edge technology should host the deception, and what
does it cost in latency, CPU, capacity and correctness?

\section{Research Questions}
\label{sec:AB.rq}
\textbf{RQ1} How much latency does WADM add per request on each edge, and how does that change with
load? \textbf{RQ2} What does it cost in CPU and capacity, and where does that cost arise?
\textbf{RQ3} Does it break what the server did without it, and does it detect what it should?
\textbf{RQ4} How much of the answer depends on the extension runtime (LuaJIT, PUC Lua,
WebAssembly)?

\section{Contributions}
(i) WADM: one configuration, five honeytoken kinds and a fake SQL-injection trap, implemented on
four edges; (ii) behavioural and tuning levelling, with an automated parity check; (iii) a paired,
three-plane benchmark harness; (iv) measured costs and where they come from.

\section{Thesis Outline}
One line per chapter, noting that Chapter~\ref{chap:AB.results} answers RQ1--RQ4 in order.

%======================================================================
\chapter{Background}
The concepts a reader needs before the design makes sense.

\section{Cyber Deception, Honeypots and Honeytokens}
Spitzner's definition \cite{spitzner2002honeypots,spitzner2003honeytokens}; honeytokens as data with
no legitimate use; why false positives are rare.

\section{Web Application Reconnaissance}
What attackers read and tamper with: HTML source, response headers, cookies, \texttt{robots.txt},
hidden form fields, SQL-injection probes.

\section{Reverse Proxies and Their Request Phases}
Request hooks, header and body filters, body buffering, upstream connection pools.

%======================================================================
\chapter{Related Work}
\label{chap:AB.related}
Prior work, grouped by theme, ending with where this thesis sits.

\section{Honeytokens and Decoys}
Decoy documents \cite{bowen2009baiting}, honeywords \cite{juels2013honeywords}, and deceptive web
artefacts such as HTML comments, \texttt{robots.txt} entries and invisible links
\cite{virvilis2014changing}.

\section{Application-Layer Deception without Code Changes}
Kahlhofer and Rass review 19 techniques for deploying deception without developer involvement,
reverse proxies among them \cite{kahlhofer2024application}.

\section{Proxy and Service-Mesh Overhead}
Sidecar overhead studies \cite{zhu2022dissecting}; Lua versus WebAssembly filters in Envoy.

\section{Benchmarking Methodology in Systems Security}
Common benchmarking flaws \cite{vanderkouwe2019sok,heiser_crimes}.

%======================================================================
\chapter{Design of WADM}
The mechanism on a conceptual level, independent of any one edge.

\section{Requirements and Threat Model}
The operator controls the edge but not the application. The attacker is a reconnaissance or
scanning client. Non-goals: blocking traffic, real authentication.

\section{Architecture}
Client $\rightarrow$ edge $\rightarrow$ static origin. One shared \texttt{config.json}, an
in-memory attacker-IP store, and the alert log as WADM's only output.
\material{\repo{CONTEXT.md} (architecture diagram)}

\section{Honeytoken Kinds}
\subsection{HTML Comments}
Injected before \texttt{</body>}; detected when the trigger keyword appears in the query.
\subsection{HTTP Headers}
A decoy response header; detected when its value is replayed in a later request.
\subsection{Cookies}
\texttt{Set-Cookie} bait; detected on tamper, since a browser replays the value unchanged.
\subsection{Decoy Paths}
Advertised as a hidden link or in \texttt{robots.txt}; any request to the path matches.
\subsection{Form Fields}
A hidden \texttt{<input>} before \texttt{</form>}; detected on tamper.

\section{The Fake SQL-Injection Trap}
The edge answers \texttt{POST /api/login} itself: a fake MySQL error that reflects the payload on a
signature hit, a canned denial otherwise. Signatures are literal strings, not regular expressions.

\section{Detection and Response Pipeline}
Parse once, scan, strip the parameter so the origin never sees it, record the IP, then write one
alert with CR/LF neutralised against log forgery.

%======================================================================
\chapter{Implementation on Four Edges}
How the design maps onto each proxy, and what made the four comparable.

\section{Common Structure}
A table mapping each edge's request hook, header filter and body filter, so the subsections below
need only describe the differences.

\section{OpenResty (Reference Edge)}
Configuration compiled once in \texttt{init\_by\_lua}, detection in \texttt{access\_by\_lua},
injection in the header and body filters, attacker IPs in a \texttt{lua\_shared\_dict}.
\material{\repo{nginx/README.md}, \repo{nginx/nginx.conf}}

\section{Envoy with a Lua Filter}
\texttt{injection.lua}, handing the request path to the response phase via dynamic metadata, body
buffering.
\material{\repo{envoy_scripts/README.md}}

\section{Envoy with a WebAssembly Filter}
Rust proxy-wasm: configuration compiled in \texttt{on\_configure}, fewer host calls, a byte-level
splice. A proxy-wasm filter cannot answer locally from the body callback.
\material{\repo{wasm-filter/README.md}, \repo{wasm-filter/src/lib.rs}}

\section{Apache httpd with mod\_lua}
The shared \texttt{wadm.lua} core and its three hook scripts; \texttt{LuaScope} and
\texttt{LuaCodeCache}.
\material{\repo{apache_scripts/README.md}}

\section{Levelling the Edges}
\label{sec:AB.level}
Product defaults levelled before any comparison: upstream keep-alive, access logs, code caching.
Baselines are the WADM configuration minus WADM. The residual differences are listed in a table.
\material{\repo{docs/EDGE_LEVELING.md}, \repo{docs/THESIS_NOTES.md} (2026-09-17, 2026-09-18)}

\section{Verifying Equivalence}
\texttt{parity\_check.py} sends 13 requests to every edge and diffs the responses, the alert lines
and the requests that reached the origin.
\material{\repo{benchmarks/parity_check.py}}

%======================================================================
\chapter{Evaluation Methodology}
How the costs were measured, and why the measurements can be trusted.

\section{Goals and Metrics}
\label{sec:AB.metrics}
A table mapping each metric to its question. RQ1: latency added (median and p95). RQ2: CPU per
request, capacity. RQ3: share of intact responses, alert recall, false positives. RQ4: in-edge
timer cost per kind.

\section{Testbed}
The VM, the cpuset layout, guest tuning, image versions.
\material{\repo{docs/BENCHMARK_METHODOLOGY.md} \S4}

\section{Workload}
k6 iterations of nine requests: a priming \texttt{GET}, four \texttt{GET} probes and four SQLi
\texttt{POST}s. Williams-balanced order \cite{williams1949designs}; arrival rates of 1, 10 and 100
iterations/s.
\material{\repo{benchmarks/test.js}}

\section{Measurement Planes}
In-edge timers (\textmu s), end-to-end latency from k6 (ms), CPU per request from cgroup
\texttt{cpu.stat}.

\section{Run Validity}
\label{sec:AB.validity}
The \texttt{assess\_run()} gates, the origin-CPU negative control, and detection of host freezes.
\material{\repo{benchmarks/wadm_timings.py}, \repo{docs/THESIS_NOTES.md} (2026-10-07)}

%======================================================================
\chapter{Results}
\label{chap:AB.results}
One section per research question, each ending with its answer.

\section{RQ1: Latency Added}
\subsection{Median Latency per GET across Rates}
$+109$ to $+218$\us{} per \texttt{GET} (13--23\%), resolved in every \texttt{GET} cell.
\material{\repo{results/paired_headline.md}, \repo{plots/paired_latency_added.png}}
\subsection{Tail Latency}
p95 grows more than the median on every edge. Apache at 100\,iterations/s collapses.
\subsection{Request Position and Reconnects}
The first request after the think-time sleep pays a wake-up cost. Apache closes the connection
after the trap's 500.
\material{\repo{docs/THESIS_NOTES.md} (2026-09-21, \S5b), \repo{plots/paired_latency_shift.png}}
\subsection{Answer to RQ1}
One paragraph: the number per edge, how it changes with load, and how certain it is.

\section{RQ2: CPU and Capacity}
\subsection{CPU per Request}
CPU added per request: OpenResty $+33$, Envoy+Lua $+41$, WASM $+123$, Apache $+200$\us.
\subsection{Cost Attribution}
The in-edge timers cover 5--22\% of that CPU; the rest is VM entry and exit, body filters and log
writes.
\material{\repo{plots/cost_attribution.png}}
\subsection{The Cost of Measuring}
The timing log lines alone cost about as much as WADM itself on OpenResty and Envoy+Lua.
\subsection{Per-Mechanism Cost and Page Size}
The \texttt{GET} path against the trap; injection cost grows linearly with body size up to
768\,KiB.
\material{\repo{results/mechanisms.md}, \repo{plots/mechanism_cpu.png}, \repo{plots/size_sweep.png}}
\subsection{Answer to RQ2}
CPU and capacity per edge, and where the cost arises.

\section{RQ3: Functional Correctness}
\subsection{Response Integrity under Load}
Share of intact responses per probe and rate. Apache is the only edge that falls below 100\%.
\material{\repo{results/paired_functional.md}}
\subsection{Detection Recall and False Positives}
\texttt{WADM ALERT} and \texttt{WADM TRAP} lines counted against expected; zero on bare.
\subsection{Answer to RQ3}
What breaks, on which edge, and whether detection is complete.

\section{RQ4: Influence of the Extension Runtime}
\subsection{Per-Kind Operation Cost}
In-edge timer distributions per honeytoken kind and edge.
\material{\repo{plots/paired_internal_timers.png}}
\subsection{The SQLi Trap: Scan Depth versus Decoding}
In the encoding-by-outcome 2$\times$2, decoding costs 0.5--2.5\us{} on the Lua edges and nothing
on WASM; scan depth is visible only on Apache's interpreted PUC Lua.
\material{\repo{results/sqli.md}, \repo{docs/THESIS_NOTES.md} (2026-09-21)}
\subsection{Stability of Apache mod\_lua}
Segfaults driven by cumulative requests in the body filter, and the cost of fixing them with
\texttt{LuaScope conn}.
\material{\repo{docs/BENCHMARK_METHODOLOGY.md} \S4b--4c, \repo{plots/apache_remediation.png}}
\subsection{Answer to RQ4}
How much of RQ1--RQ3 is explained by the runtime rather than by the deception logic.

%======================================================================
\chapter{Discussion}
What the results mean, and how far they can be trusted.

\section{Synthesis across Research Questions}
One table of edges against the four answers.

\section{Choosing an Edge for Deception}
A trade-off summary per edge: latency, CPU, capacity, robustness.

\section{Where the Cost Comes From}
Runtime entry and exit, body buffering and log writes, not the detection logic.

\section{Deployment Implications}
Compressed bodies, range requests, preserved ETags, body-size limits.

\section{Threats to Validity}
A single VM host, synthetic traffic, moving image tags, timer resolution.

\section{Limitations}
POST-body inspection exists on only two edges; no real attacker traffic; no TLS or HTTP/2.

%======================================================================
\chapter{Conclusion and Future Work}
\section{Summary}
RQ1--RQ4 answered in one or two sentences each, then the overall contribution.
\section{Future Work}
POST-body inspection on all four edges, injecting into compressed bodies, a field study with real
attackers, TLS and HTTP/2.

%======================================================================
\appendix
\chapter{Configuration Reference}
Every \texttt{config.json} field per honeytoken kind, and the SQLi trap's fields.
\material{\repo{CONTEXT.md}}
\chapter{Parity-Check Request Set}
The 13 requests and the documented exclusions.
\chapter{Reproducing the Benchmarks}
Commands, machine configuration, run durations.
\material{\repo{docs/BENCHMARK_METHODOLOGY.md} \S3--4}
\chapter{Additional Result Tables}
Per-rate tables and figures.
\material{\repo{benchmarks/results/plots/by_rate/}}

\bibliographystyle{ieeetr}
\bibliography{sources}
\end{document}
